\documentclass[12pt,a4paper]{article}

% Packages
\usepackage{amsmath, amssymb, amsthm}
\usepackage{booktabs}
\usepackage[table]{xcolor}
\usepackage{geometry}
\usepackage{hyperref}
\usepackage{enumitem}
\usepackage{mathtools}
\usepackage{fancyhdr}
\usepackage{tcolorbox}

% Page geometry
\geometry{margin=1in}

% Hyperlink setup
\hypersetup{
    colorlinks=true,
    linkcolor=blue!70!black,
    urlcolor=cyan!70!black
}

% Commands
\newcommand{\QSVT}{\textsc{qsvt}}
\newcommand{\LCU}{\textsc{lcu}}
\newcommand{\QFT}{\textsc{qft}}
\newcommand{\poly}{\operatorname{poly}}
\newcommand{\polylog}{\operatorname{polylog}}
\newcommand{\cyc}{\mathbb{Q}(\zeta_n)}

% Header/Footer
\pagestyle{fancy}
\fancyhf{}
\rhead{\thepage}
\lhead{Algebraic Structure of Spectral Algorithms \quad June 2026}

\begin{document}

\title{
\textbf{
Strategy Proposal:\\
Algebraic Field Tracking in Spectral Quantum Algorithms
}
}

\author{
Exploratory Framework for Algebraically Structured FTQC
}

\date{June 2026}

\maketitle

\begin{abstract}
This document outlines a speculative research program investigating the algebraic number-theoretic structure underlying spectral quantum algorithms based on block encodings, Chebyshev expansions, QSVT, LCU, and qubitization. 

Rather than optimizing query complexity or T-count directly, the goal is to classify the algebraic field extensions generated throughout a quantum spectral-processing pipeline. Preliminary observations suggest that restricted-phase Chebyshev-QSVT architectures may generate remarkably low-complexity algebraic structures, with non-cyclotomic extensions appearing primarily through global normalization steps.

The long-term objective is to explore whether algebraic field complexity can serve as an alternative organizational principle for fault-tolerant compilation, exact synthesis, and future quantum gate architectures.
\end{abstract}

\tableofcontents

\vspace{1em}
\hrule
\vspace{1em}

% ========================================================

\section{Motivation}

Fault-tolerant quantum compilation is traditionally analyzed through:

\begin{itemize}
    \item T-count,
    \item magic-state consumption,
    \item circuit depth,
    \item approximation complexity.
\end{itemize}

However, many synthesis methods are fundamentally controlled by algebraic number theory.

Examples include:

\begin{itemize}
    \item Clifford+T exact synthesis,
    \item cyclotomic unitary representations,
    \item metaplectic compiling,
    \item ancilla-assisted algebraic synthesis.
\end{itemize}

This motivates the following broader question:

\begin{quote}
What algebraic field extensions are naturally generated by spectral quantum algorithms?
\end{quote}

The present strategy proposes a systematic investigation of this question in the setting of:

\begin{itemize}
    \item block encodings,
    \item Chebyshev-QSVT,
    \item LCU constructions,
    \item qubitization/amplitude amplification,
    \item iterative spectral transforms.
\end{itemize}

% ========================================================

\section{Guiding Hypothesis}

Preliminary observations suggest the following principle:

\begin{quote}
Many structured spectral quantum algorithms generate surprisingly low-degree algebraic field extensions despite operating on highly nontrivial operators.
\end{quote}

In particular:

\begin{enumerate}
    \item Chebyshev-QSVT layers often require only small cyclotomic phase sets.
    
    \item Polynomial coefficients generated through iterative schemes (e.g.\ Newton--Schulz) may remain rational.
    
    \item The dominant algebraic extensions frequently appear through global normalization factors such as:
    \[
    \frac1{\sqrt{\lambda}}.
    \]
    
    \item These normalization steps often introduce only quadratic extensions.
\end{enumerate}

The resulting architecture may therefore live inside very restricted towers of algebraic number fields.

% ========================================================

\section{Primary Research Questions}

The project investigates the following questions.

\subsection{Field Growth}

Given a spectral algorithm pipeline:
\[
A
\rightarrow
p(A)
\rightarrow
\QSVT
\rightarrow
\LCU
\rightarrow
\text{amplification},
\]
what is the minimal algebraic field containing all amplitudes and phases?

\subsection{Extension Degree}

Does the extension degree grow:

\begin{itemize}
    \item logarithmically,
    \item polynomially,
    \item exponentially,
    \item or remain bounded?
\end{itemize}

\subsection{Quadratic Dominance}

To what extent are nontrivial extensions generated only by iterated quadratic radicals?

\subsection{Cyclotomic Restrictions}

Can restricted-phase QSVT architectures be classified by their underlying cyclotomic fields?

% ========================================================

\section{Reference Example: Matrix Inversion Pipeline}

The proposed matrix inversion strategy serves as the initial case study.

\subsection{Newton--Schulz Layer}

The iteration:
\[
X_{k+1}
=
X_k(2I-AX_k)
\]
preserves rationality whenever:
\[
A \in \mathbb Q^{N\times N}
\]
and the initialization is rational.

Thus:
\[
p_k(x)\in\mathbb Q[x].
\]

\subsection{Chebyshev Expansion}

The basis conversion:
\[
p_k(x)
=
\sum_m c_m T_m(x)
\]
also preserves rationality:
\[
c_m\in\mathbb Q.
\]

\subsection{Restricted-Phase QSVT}

Chebyshev iterates admit restricted phase sequences often chosen from:
\[
\left\{
0,\frac{\pi}{2},\pi,-\frac{\pi}{2}
\right\}.
\]

Consequently, the corresponding gates lie inside small cyclotomic fields:
\[
\mathbb Q(i),
\qquad
\mathbb Q(\zeta_8),
\]
or nearby extensions.

\subsection{Global Normalization}

The first major non-cyclotomic contribution arises from:
\[
\lambda
=
\sum_m |c_m|,
\]
through normalization factors:
\[
\frac1{\sqrt{\lambda}}.
\]

This introduces at most a quadratic extension:
\[
\mathbb Q(\sqrt{\lambda}).
\]

% ========================================================

\section{Proposed Algebraic Complexity Measure}

The project proposes studying alternatives to standard resource measures.

Possible candidates include:

\begin{enumerate}
    \item extension degree,
    \item cyclotomic conductor,
    \item radical depth,
    \item Galois group complexity,
    \item tower height of quadratic extensions.
\end{enumerate}

One possible measure is:

\[
\mathcal C_{\mathrm{alg}}
=
[\mathbb K : \mathbb Q],
\]
where:
\[
\mathbb K
\]
is the smallest field containing all circuit amplitudes.

Another possibility is tracking only the non-cyclotomic component:
\[
[\mathbb K : \cyc].
\]

% ========================================================

\section{Potential Connections}

The framework is naturally connected to several existing areas.

\subsection{Exact Quantum Synthesis}

Connections include:

\begin{itemize}
    \item Kliuchnikov--Maslov--Mosca exact synthesis,
    \item cyclotomic unitary representations,
    \item Clifford+T algebraic normal forms,
    \item Ross--Selinger synthesis,
    \item GridSynth methods.
\end{itemize}

\subsection{Alternative Gate Sets}

Restricted algebraic growth may become highly relevant for:

\begin{itemize}
    \item metaplectic gate models,
    \item qudit architectures,
    \item algebraically adapted FTQC designs,
    \item exact ancilla-assisted synthesis.
\end{itemize}

\subsection{Compiler Theory}

Future compilers may potentially optimize:

\begin{itemize}
    \item algebraic field degree,
    \item exact representability,
    \item radical structure,
    \item cyclotomic embeddings,
\end{itemize}

rather than only T-count.

% ========================================================

\section{Proposed Technical Program}

The project may proceed through several stages.

\subsection{Stage 1: Symbolic Field Tracking}

For specific spectral algorithms:

\begin{itemize}
    \item compute all generated amplitudes symbolically,
    \item determine minimal enclosing fields,
    \item classify extension towers.
\end{itemize}

\subsection{Stage 2: Restricted QSVT Classification}

Study:

\begin{itemize}
    \item which polynomial families admit low-field QSVT realizations,
    \item relations between phase structure and field growth,
    \item Clifford-dominant versus arbitrary-phase QSVT.
\end{itemize}

\subsection{Stage 3: Exact Synthesis Implications}

Investigate:

\begin{itemize}
    \item whether bounded algebraic complexity yields cheaper synthesis,
    \item specialized compiler strategies,
    \item ancilla-assisted realizations,
    \item alternative FT architectures.
\end{itemize}

\subsection{Stage 4: Structured Physics Applications}

Potential targets include:

\begin{itemize}
    \item lattice field theory,
    \item sparse PDE discretizations,
    \item translation-invariant Hamiltonians,
    \item quantum walk architectures,
    \item structured spectral filters.
\end{itemize}

% ========================================================

\section{Relation to the Low-T Program}

This project should not be viewed as replacing the low-T Chebyshev-QSVT framework.

Instead:

\begin{itemize}
    \item the low-T program studies compiler resource localization,
    \item the present program studies algebraic field localization.
\end{itemize}

The two perspectives are closely related but conceptually distinct.

Low algebraic complexity does not automatically imply low T-count, but may nevertheless reveal deep structural constraints relevant for future compilation paradigms.

% ========================================================

\section{Long-Term Vision}

The long-term hypothesis is:

\begin{quote}
Spectral quantum algorithms may admit a classification theory based on algebraic field structure, analogous to how classical harmonic analysis classifies transforms through symmetry and representation theory.
\end{quote}

In this picture:

\begin{itemize}
    \item arbitrary-phase QSVT represents the fully general setting,
    \item restricted Chebyshev-QSVT occupies a low-field subclass,
    \item and structured physical operators may naturally generate bounded algebraic complexity.
\end{itemize}

Such a classification could eventually influence:

\begin{itemize}
    \item fault-tolerant compilation,
    \item exact synthesis,
    \item hardware gate design,
    \item algebraically aware compilers,
    \item and structured spectral algorithm design.
\end{itemize}

% ========================================================

\section{Conclusion}

We proposed a speculative research strategy investigating the algebraic field structure generated by spectral quantum algorithms.

The primary conjectural insight is that restricted-phase spectral architectures may generate surprisingly small algebraic extensions, with most nontrivial field growth arising from global normalization operations.

Rather than focusing solely on asymptotic query complexity or T-count, the project introduces the possibility of:

\[
\textit{algebraic field complexity}
\]
as an additional organizational principle for fault-tolerant quantum computation.

\vspace{2em}

\hrule

\vspace{1em}

\noindent
\textit{
Exploratory strategy document generated from collaborative discussion.
}

\end{document}